\documentclass[10pt]{article}
\usepackage{icmlko}

\icmlmeta{IP 파운데이션 월드모델}{90}{네 번째 축을 네 번 시도했다}{넷 다 떨어졌고, 소거법으로 재니 현행 설계가 국소 최적이었다 --- 그리고 분모가 또 속일 뻔했다}

\begin{document}

\twocolumn[
\icmlheader

\begin{abstract}
노트 89가 숙제를 남겼다 --- 기획서 열 칸 중 모델이 버리는 다섯 축의 대응물을
다른 도메인에서 찾아라. \textbf{체험 밀도}를 열 도메인 전부에서 만들었다(게임
기능 수 · 도서 쪽수 · 펀딩 리워드 단계 수 · 웹툰/애니/만화/세계애니 회차 수 ·
모바일 스크린샷 수 · 팝업 사람 태그). \textbf{떨어졌다} --- 공통 축으로
$-$0.060, 고유 축으로 $-$0.029다. 이것으로 \textbf{네 번째 공통 축 시도가
네 번 모두 실패}했다: 경쟁 밀도 $-$0.020(노트 84) · 전형성 $-$0.009(84) ·
IP 인지도 팝업 $-$0.051(89) · 체험 밀도 $-$0.060. 그래서 축을 더하는 대신
\textbf{빼면서} 설계를 검사했다. 공통 축 셋 중 타깃 폭을 빼면 $-$0.028(팝업은
$+$0.450에서 $+$0.170으로 무너진다), 굿즈 규모를 빼면 $-$0.028이다. 성분 수는
$k{=}2$가 $k{=}1$보다 $+$0.019, $k{=}3$보다 $+$0.006 낫다. \textbf{현행
설계(공통 3축 · 성분 2)가 국소 최적이다.} 다만 하나가 갈린다 --- \textbf{팝업만
보면 $k{=}1$이 $+$0.4736으로 현행($+$0.4501)보다 낫다.} 그리고 함정을 하나
또 밟을 뻔했다. 매장 노출도를 빼면 판정치가 $+$0.0077 \emph{오르는 것처럼}
보이는데, 그 축이 없으면 정렬이 안 되는 도메인이 생겨 \textbf{대상이 10개에서
9개로 준} 것이었다. 고정하고 재면 $-$0.0067이다 --- \textbf{노트 82의 분모
문제가 축 수준에서 그대로 재현됐다.}
\end{abstract}
]

\begin{figure*}[t]
\centering
\includegraphics[width=\textwidth]{figs/fourth.pdf}
\caption{왼쪽 --- 네 번째 공통 축 시도 넷. 오른쪽 --- 공통 축 수와 성분 수의 격자.}
\end{figure*}

\section{체험 밀도 --- 네 번째 실패}

``안에서 할 일이 얼마나 되나''는 어느 도메인에나 있다. 열 도메인 전부에서
만들었다.

\begin{table}[t]
\centering\tsize
\caption{체험 밀도의 대응물.}
\begin{tabular}{ll}
\toprule
팝업 & 체험 밀도(사람이 매김) \\
게임 & 기능 수 \\
도서 & 쪽수 \\
펀딩 & 리워드 단계 수 \\
웹툰 · 애니 · 만화 · 세계애니 & 회차 수 \\
모바일 & 스크린샷 수 \\
\bottomrule
\end{tabular}
\end{table}

\begin{table}[t]
\centering\tsize
\caption{판정치와 팝업. 대상 10개 고정.}
\begin{tabular}{lrr}
\toprule
& 판정치 & 팝업 \\
\midrule
현행 & \textbf{$+$0.4038} & \textbf{$+$0.4501} \\
공통 4축(중립 대입) & $+$0.3439 & $+$0.4377 \\
고유 축(중립 대입) & $+$0.3745 & $+$0.4240 \\
공통 4축(마스크 유지) & $+$0.3597 & $+$0.4333 \\
\bottomrule
\end{tabular}
\end{table}

\begin{table}[t]
\centering\tsize
\caption{네 번째 공통 축 시도 넷.}
\begin{tabular}{lrr}
\toprule
& $\Delta$판정치 & $\Delta$팝업 \\
\midrule
경쟁 밀도(노트 84) & $-$0.020 & --- \\
전형성(노트 84) & $-$0.009 & --- \\
IP 인지도(노트 89) & $+$0.011 & \textbf{$-$0.051} \\
\textbf{체험 밀도(노트 90)} & \textbf{$-$0.060} & $-$0.012 \\
\bottomrule
\end{tabular}
\end{table}

\claimx{공통 축을 넷째로 늘리는 것은 닫혔다 --- 네 번 시도해 네 번 떨어졌다.
성질이 다른 넷(시장 혼잡도 · 내용 전형성 · 사전 인지도 · 내용 분량)이 전부
같은 결과를 냈다.}{넷 다 ``도메인마다 다른 물리량을 같은 슬롯에 넣는'' 방식이었다.
열 도메인이 문자 그대로 같은 것을 재는 양은 아직 경쟁 밀도 하나뿐이었고
그것도 시간과 겹쳐 실패했다.}

\section{빼면서 검사한다}

더하는 것이 안 되니 빼 봤다.

\begin{table}[t]
\centering\tsize
\caption{공통 축을 하나씩 뺄 때. \textbf{대상 9개 고정}.}
\begin{tabular}{lrr}
\toprule
& 판정치 & 팝업 \\
\midrule
\textbf{현행 3축} & \textbf{$+$0.4182} & $+$0.4501 \\
매장 노출도 뺌 & $+$0.4115 & $+$0.4499 \\
\midrule
\multicolumn{3}{l}{\emph{(대상 고정 안 하면)} 매장 노출도 뺌 $+$0.4115 $\to$
``$+$0.0077 상승''} \\
\bottomrule
\end{tabular}
\end{table}

\claimx{매장 노출도를 빼면 판정치가 오르는 것처럼 보이지만 분모 효과다. 그 축이
없으면 정렬이 안 되는 도메인이 생겨 대상이 10개에서 9개로 줄고, 빠지는 것이
성적이 낮은 대상이다. 고정하고 재면 $-$0.0067이다.}{9개로 고정한 비교도
완전하지는 않다 --- 빠진 대상(애니)이 어느 설정에서 살아나느냐가 결과를 바꿀
수 있다.}

\textbf{노트 82가 도메인 수준에서 발견한 함정이 축 수준에서 그대로 재현됐다.}
그때 세운 규약 --- ``분모가 바뀌는 비교에는 다른 자가 필요하다'' --- 이 여기서도
그대로 적용된다.

\section{성분 수}

\begin{table}[t]
\centering\tsize
\caption{공통 축 수 $\times$ 성분 수. 대상 9개 고정.}
\begin{tabular}{lrrr}
\toprule
& $k{=}1$ & $k{=}2$ & $k{=}3$ \\
\midrule
공통 3축 (평균) & $+$0.3990 & \textbf{$+$0.4182} & $+$0.4121 \\
공통 3축 (팝업) & \textbf{$+$0.4736} & $+$0.4501 & $+$0.4216 \\
\midrule
공통 2축 (평균) & $+$0.3990 & $+$0.4115 & $+$0.4112 \\
공통 2축 (팝업) & $+$0.4736 & $+$0.4499 & $+$0.4464 \\
\bottomrule
\end{tabular}
\end{table}

$k{=}1$에서는 공통 축 셋과 둘이 \emph{똑같다} --- 성분이 하나면 셋째 축이
회전에 기여할 자리가 없기 때문이다. 노트 37이 ``공통 축 $c$, 성분 $k$일 때
$c<k$면 회전이 과소결정된다''고 했는데, 반대쪽 경계도 있다 --- \textbf{$k{=}1$이면
$c$가 몇이든 같다.}

\claimx{현행 설계(공통 3축 · 성분 2)가 국소 최적이다. 축을 더해도 빼도, 성분을
늘려도 줄여도 판정치가 내린다.}{$k{=}1$이 팝업에서는 $+$0.4736으로 현행보다
$+$0.0235 낫다. 노트 75가 보인 대로 팝업 75건으로는 그 차이를 잡음과 못
가르지만, \textbf{제품 목표만 보면 다른 설계가 있을 수 있다는 신호다.}}

\section{한계}

\begin{itemize}[leftmargin=1.1em,itemsep=0.15em]
\item 체험 밀도의 대응물이 도메인마다 성질이 다르다 --- 회차 수는 누적이고
      스크린샷 수는 마케팅 투입이다. ``할 일의 양''이라는 하나의 뜻으로
      묶은 것이 무리였을 수 있다.
\item 남은 세 축(포토존 · 컬래버 강도 · 계절 적합)은 안 만들었다. 대응물이
      더 얇다고 보고 멈췄다.
\item 소거법을 공통 축 하나씩만 했다. 둘을 동시에 빼면 다를 수 있다.
\item $k{=}1$이 팝업에서 낫다는 것을 채택 검정으로 확인하지 않았다.
\item 대상 9개 고정 비교에서 애니가 빠졌다 --- 애니는 굿즈 규모가 꺼져 있어
      축을 빼면 정렬이 먼저 무너지는 도메인이다.
\end{itemize}

\section{다음}

\begin{enumerate}[leftmargin=1.3em,itemsep=0.15em]
\item \textbf{제품 전용 설계를 따로 둘지 정한다} --- 팝업만 보면 $k{=}1$이
      $+$0.0235 낫다. 판정치용과 제품용을 갈라 두는 것이 정직한지 검토한다.
\item 열 도메인이 \emph{문자 그대로} 같은 것을 재는 양을 더 찾는다. 경쟁 밀도가
      유일했고 시간과 겹쳐 실패했다.
\item 제품 도구(노트 89)를 다른 카테고리로 넓힌다 --- 웹툰 · 만화 신작 기획.
\item 팝업 표본. 네 노트가 같은 곳을 가리킨다(75 · 86 · 87 · 90).
\end{enumerate}

\vspace{0.6em}
\hrule height 0.3pt
\vspace{0.5em}
{\footnotesize
\textbf{참고문헌}\\[0.3em]
Schoenemann, P. H. (1966). A generalized solution of the orthogonal Procrustes
problem. \emph{Psychometrika}, 31(1), 1--10.\\[0.15em]
Meredith, W. (1993). Measurement invariance, factor analysis and factorial
invariance. \emph{Psychometrika}, 58(4), 525--543.\\[0.15em]
Simpson, E. H. (1951). The interpretation of interaction in contingency tables.
\emph{JRSS-B}, 13(2), 238--241.\\[0.15em]
Ben-David, S. et al. (2010). A theory of learning from different domains.
\emph{Machine Learning}, 79(1--2), 151--175.\\[0.15em]
Efron, B., Tibshirani, R. J. (1993). \emph{An Introduction to the Bootstrap}.
Chapman \& Hall.
}

\end{document}
